%% ================================================================
%% SFST v30 — Appendix V: Equipartition of SU(3) gauge fluctuations
%% 
%% Tier structure:
%%   Theorem V.1 (Schur): Tier 1 — rigorous, no assumptions beyond Lie theory
%%   Corollary V.2 (Equipartition): Tier 1 — given Tier 1 hypotheses
%%   Proposition V.3 (σ/√8): Tier 2 — conditional on SU(3)-invariance of vacuum
%%   Remark V.4 (σ=1): Tier 3 — normalization not yet derived
%%
%% Key result: The 1/√8 factor is not numerology — it follows from
%% Schur's Lemma applied to the irreducible adjoint representation of SU(3).
%% The remaining open point is the normalization σ=1, not the √8 structure.
%% ================================================================

\section{Appendix V: Equipartition of $\mathrm{SU}(3)$ gauge fluctuations
  and the algebraic origin of $1/\sqrt{8}$}
\label{app:equipartition}

This appendix supplies the mathematical proof that the factor $1/\sqrt{8}$
in the SFST formula follows from Schur's Lemma applied to the irreducible
adjoint representation of $\mathrm{SU}(3)$.  The proof clarifies which
parts of the argument are Tier~1 (no assumptions) and which require an
additional normalization condition.

%% ----------------------------------------------------------------
\subsection{Setup and statement}

Let $A^a_\mu$ ($a=1,\ldots,8$) be the components of an $\mathrm{SU}(3)$
gauge field in the adjoint representation, with vacuum two-point function
\[
  M^{ab} \;\equiv\; \bigl\langle A^a_\mu A^b_\mu \bigr\rangle_{\mathrm{vac}}.
\]
Denote by $\sigma^2 \equiv \sum_{a=1}^8 M^{aa}$ the total fluctuation energy,
assumed finite (see Remark~\ref{rem:finite}).

\begin{theorem}[Equipartition, Tier~1]
  \label{thm:equipartition}
  Let hypotheses \textup{(H1)--(H3)} hold:
  \begin{enumerate}[label=\textup{(H\arabic*)}]
    \item The vacuum state $\langle\cdot\rangle_{\mathrm{vac}}$ is invariant
          under $\mathrm{SU}(3)_{\mathrm{color}}$ gauge transformations.
    \item The adjoint representation
          $\mathrm{Ad}\colon\mathrm{SU}(3)\to\mathrm{GL}(\mathfrak{su}(3))$
          is irreducible.
    \item The total fluctuation energy $\sigma^2 = \sum_a M^{aa}$ is finite.
  \end{enumerate}
  Then:
  \begin{equation}
    M^{ab} = \frac{\sigma^2}{8}\,\delta^{ab}
    \qquad\text{for all }a,b=1,\ldots,8,
  \end{equation}
  and in particular $\langle(A^a)^2\rangle = \sigma^2/8$ for every
  generator $a$.  The amplitude per channel is $\sigma/\sqrt{8}$.
\end{theorem}

\begin{proof}
By~(H1), the two-point function $M^{ab}$ is equivariant under the adjoint action:
\[
  M^{ab} = \bigl\langle A^a A^b\bigr\rangle
    = \bigl\langle (\mathrm{Ad}_g)^{a}{}_{c}\,A^c\;
                  (\mathrm{Ad}_g)^{b}{}_{d}\,A^d\bigr\rangle
    = (\mathrm{Ad}_g)^{a}{}_{c}\,(\mathrm{Ad}_g)^{b}{}_{d}\,M^{cd}
\]
for every $g\in\mathrm{SU}(3)$.  In other words, $M$ is an
$\mathrm{SU}(3)$-invariant bilinear form on $\mathfrak{su}(3)$.

By~(H2), $\mathrm{Ad}$ is irreducible.  Schur's Lemma then implies that every
$\mathrm{SU}(3)$-equivariant endomorphism of $\mathfrak{su}(3)$ is a scalar
multiple of the identity.  Since $M$ defines such an endomorphism via
$X\mapsto M(X,\cdot)$, we have $M^{ab}=\lambda\,\delta^{ab}$ for some
$\lambda\in\mathbb{R}$.

By~(H3), $\sum_{a=1}^8 M^{aa} = 8\lambda = \sigma^2$, giving
$\lambda=\sigma^2/8$.
\end{proof}

%% ----------------------------------------------------------------
\subsection{Why hypothesis (H2) holds: the adjoint of $\mathrm{SU}(3)$
  is irreducible}

\begin{proof}[Proof of (H2)]
The Lie algebra $\mathfrak{su}(3)$ is \emph{simple}: its only ideals are
$\{0\}$ and $\mathfrak{su}(3)$ itself.  A subrepresentation of the adjoint
representation is precisely an ideal of $\mathfrak{su}(3)$.  Since
$\mathfrak{su}(3)$ is simple, the adjoint representation admits no
non-trivial subrepresentations and is therefore irreducible.  $\square$

In Dynkin notation, the adjoint of $\mathrm{SU}(3)$ carries highest weight
$(1,1)$ and dimension $8 = N_c^2-1\big|_{N_c=3}$.
\end{proof}

%% ----------------------------------------------------------------
\subsection{The connection to $1/\sqrt{8}$}

\Cref{thm:equipartition} gives the per-channel amplitude as $\sigma/\sqrt{8}$,
where $\sigma^2$ is the total fluctuation energy.  The specific value
$1/\sqrt{8}$ in the SFST formula requires $\sigma=1$.

\begin{proposition}[Conditional normalisation, Tier~2]
  \label{prop:sigma1}
  Suppose, in addition to \textup{(H1)--(H3)}, that the SFST saturation
  condition fixes $\sigma^2=1$ in the natural spectral units of
  $V_{\mathrm{eff}}(R_0)$.  Then the per-channel amplitude equals exactly
  $1/\sqrt{8}$, and the collective electromagnetic correction takes the form
  $\alpha^2/\sqrt{8}$.
\end{proposition}

The condition $\sigma^2=1$ is not yet derived from the SFST effective
potential.  It would follow if one could show that the zeta-regularised
total gauge fluctuation $\sigma^2 = \sum_a\langle(A^a)^2\rangle_{\mathrm{SFST}}$
equals the unit spectral action value at $R_0=1/2$.  This is an open
technical problem.

%% ----------------------------------------------------------------
\subsection{What the Higgs comparison establishes}

The claim that equipartition implies SFST operates ``below symmetry
breaking'' requires care.  Three symmetry-breaking scales are relevant
in the Standard Model:
\begin{enumerate}[label=(\alph*)]
  \item \textbf{Electroweak breaking}
    ($\mathrm{SU}(2)_L\times\mathrm{U}(1)_Y\to\mathrm{U}(1)_{\mathrm{em}}$).
    This affects $\mathrm{SU}(2)_L$ and the electroweak $\mathrm{U}(1)$.
    The $\mathrm{SU}(3)_{\mathrm{color}}$ gauge symmetry is
    \emph{not} broken by the electroweak Higgs.
    SFST couples to $\mathrm{SU}(3)_{\mathrm{color}}$, so electroweak
    breaking is irrelevant.
  \item \textbf{QCD confinement.}  Confinement does \emph{not} break
    $\mathrm{SU}(3)_{\mathrm{color}}$ spontaneously; there is no
    associated Goldstone boson.  The $\mathrm{SU}(3)_{\mathrm{color}}$
    symmetry of the vacuum is intact, so equipartition holds in the
    confined phase.
  \item \textbf{Chiral breaking}
    ($\mathrm{SU}(3)_L\times\mathrm{SU}(3)_R\to\mathrm{SU}(3)_V$).
    This breaks $\mathrm{SU}(3)_{\mathrm{flavor}}$, not
    $\mathrm{SU}(3)_{\mathrm{color}}$, and does not affect the
    adjoint color correlator.
\end{enumerate}
In summary: $\mathrm{SU}(3)_{\mathrm{color}}$ equipartition holds in
every physically realised phase of QCD.  The statement ``SFST operates
below symmetry breaking'' is therefore automatically satisfied in the
color sector and does not constitute an additional constraint.

The genuine content of hypothesis (H1) is weaker: it requires that the
SFST ground state at $R_0=1/2$ be $\mathrm{SU}(3)_{\mathrm{color}}$-invariant.
This is physically motivated (the T-duality minimum has no preferred colour
direction) but has not been proved from the SFST Lagrangian.

%% ----------------------------------------------------------------
\subsection{Summary and tier classification}

\begin{center}
\renewcommand{\arraystretch}{1.3}
\begin{tabular}{p{0.40\textwidth}cp{0.30\textwidth}}
\toprule
Statement & Tier & Status \\
\midrule
$\mathrm{Ad}(\mathrm{SU}(3))$ is irreducible &
  1 & proved: $\mathfrak{su}(3)$ simple \\
$M^{ab}=\lambda\delta^{ab}$ in any SU(3)-invariant state &
  1 & proved: Schur's Lemma \\
Per-channel amplitude $= \sigma/\sqrt{8}$ &
  1 & proved: (H1)--(H3) \\
SFST vacuum is SU(3)-invariant (H1) &
  2 & motivated, not proved \\
Normalisation $\sigma=1$ from saturation &
  3 & open \\
\bottomrule
\end{tabular}
\end{center}

\begin{remark}[Finiteness of $\sigma^2$]
  \label{rem:finite}
  On the flat $T^5$, the vacuum two-point function requires
  UV regularisation.  In the zeta-regularisation scheme
  used throughout this paper, $V_{\mathrm{eff}}(R_0)$ is finite
  (computed numerically in Appendix~S).  Whether
  $\sigma^2=\sum_a\langle(A^a)^2\rangle$ in the same scheme is
  equal to $V_{\mathrm{eff}}(R_0)$ or to a different finite combination
  is the technical open point noted in Proposition~\ref{prop:sigma1}.
\end{remark}

The Equipartition Theorem establishes that $1/\sqrt{8}$ is not numerology
but a direct consequence of Schur's Lemma applied to the irreducible
$\mathrm{SU}(3)$ adjoint.  The remaining open problem is not the
$\sqrt{8}$ structure but the identification of $\sigma=1$ with the
SFST saturation energy scale.
